\documentclass[10pt,a4paper]{amsart}
\usepackage{amsmath,amssymb,amsthm,graphicx}
\usepackage{mathrsfs}
\newtheorem*{definition}{Definition}
\newtheorem*{theorem}{Theorem}
\newcommand{\ES}{\mathbf{Ev}}
\newcommand{\SES}{\mathbf{EvSym}}
\newcommand{\Cat}{\mathbf{Cat}}
\newcommand{\Con}{\mathrm{Con}}
\newcommand{\conf}{\mathscr C}
\newcommand{\id}{\mathrm{id}}
\AtBeginDocument{\pagestyle{plain}}
\newcommand{\Petri}{\mathbf{Petri}}
\newcommand{\Occ}{\mathbf{Occ}}
\title{The embedding of event structures into event structures with symmetry is 2-dense}
\author{Serge Lechenne, Hugo Paquet}
\date{arxiv:2609.17324v1; CC BY4.0 (fixed arXiv2609.17324v1 native author work)}
\begin{document}
\maketitle
\noindent\textit{Editorial note.} Source https://arxiv.org/abs/2609.17324v1. Copyright Serge Lechenne and Hugo Paquet, CC BY4.0. Complete native author 2-density construction retained, including its explicit omitted-verification statement. Existing corpus210587 counts only prop:Qonmor; this item is the independent main event-structure 2-density theorem (conference Theorem36), not a duplicate or a new count for its supporting context.\par\medskip
\noindent\textit{Allowed standard prerequisites (editorial).}\begin{itemize}
\item Locally small 2-categories, enriched hom functors, strict natural transformations and modifications. Native 2-density definition and publisher reference17 Kelly; event-structure symmetry conventions from reference32 and actual native definitions included here.
\end{itemize}\medskip
\begin{definition}
An \emph{event structure} is a tuple $(E, {\leq}, \#)$ where $\leq$ is a partial order on $E$, and $\#$  is an irreflexive and symmetric binary relation on $E$, satisfying the axioms below:
\begin{itemize}
    \item For every $e \in E$  the set $\{e' \in E \mid e' \leq e \}$ is finite.
    \item For $e, e', e'' \in E$, if $e \leq_{E} e'$ and $e \mathrel{\#} e''$ then $e' \mathrel{\#} e''$.
\end{itemize}
\end{definition}
The execution paths of an event structure are called \emph{configurations}:
\begin{definition}
    A finite subset $X \subseteq E$ is called \emph{consistent} if no two events in $X$ are in conflict. It is called a \emph{configuration} if it is additionally down-closed. The set of consistent subsets of $E$ is denoted $\Con(E)$, and the set of configurations is called $\conf(E)$.
\end{definition}
Maps of event structures must faithfully preserve configurations:
\begin{definition}
    If $(E, \leq_E, \#_E)$ and $(E', \leq_{E'}, \#_{E'})$ are event structures, a function $f: E \rightarrow E'$ is a \emph{map of event structures} if for every $x \in \conf(E)$, the direct image $fx := \{f(e) \mid e \in x\}$ is in $\conf(E')$, and moreover the restriction of $f$ to $x$ is injective.
\end{definition}
Event structures and maps of event structures form a category $\ES$. Frequently, we will regard $\ES$ as a locally discrete $2$-category, just like we did for $\Petri$ and $\Occ$.

\subsection{The unfolding of a Petri net as an event structure with symmetry}

We first recall event structures with symmetry and the $2$-category $\SES$, originally introduced in [32].

\subsubsection{Event structures with symmetry}

\label{subsec:sem-def}
 Informally, symmetry on an event structure indicates which pairs of executions can be considered equivalent, via a bisimulation relation. 
\begin{definition}
    For an event structure $E$, an isomorphism family  is a set $\mathbb{S}$ of bijections $\theta:x \cong y$ between finite configurations of $x, y \in \conf(E)$, such that:
    \begin{itemize}
        \item $\mathbb{S}$ contains all identity bijections, and is stable under composition and inverse.  
        \item For $\theta : x \cong y \in \mathbb{S}$, if $x' \subseteq x$, then the restriction of $\theta$ to $x'$ is in $\mathbb{S}$.
        \item For $\theta : x \cong y \in \mathbb{S}$, if $x \subseteq x'$, then there exists an extension $y \subseteq y'$ and a bijection $\theta' : x' \cong y' \in \mathbb{S}$ such that $\theta'$ restricts to $\theta$.
    \end{itemize}
    The pair $(E, \mathbb{S})$ is called an \emph{event structure with symmetry}.
\end{definition}

We then consider symmetry-preserving maps:
\begin{definition}
    For event structures with symmetry $(E,\mathbb{S})$ and $(E', \mathbb{S}')$, a map $f : E \rightarrow E'$ \emph{preserves symmetry} if for all $\theta : x \cong y \in \mathbb{S}$, $f \theta : fx \cong fy \in \mathbb{S}'$, where $f\theta : f(e) \mapsto f(\theta(e))$.
\end{definition}
A novelty in the presence of symmetry is the following equivalence relation on maps:
\begin{definition}
    For event structures with symmetry $(E,\mathbb{S})$ and $(E', \mathbb{S}')$ and symmetry-preserving maps $f, g : E \rightarrow E'$, say that $f$ and $g$ are \emph{symmetric}, denoted $f \sim g$, if for every $x \in \conf(E)$ the bijection $fx \cong gx$ defined by $ f(e) \mapsto g(e)$ is in $\mathbb{S}'$.
\end{definition}
Together, event structures with symmetry, symmetry-preserving maps and symmetries of maps form a $2$-category $\SES$ [32]. This is a degenerate 2-category in the sense that hom-categories are all \emph{setoids} (sets with an equivalence relation). There is an embedding $J : \ES \to \SES$. As we will see, the 2-categorical structure makes it possible to consider universal properties ``up to symmetry''.

\subsection{A 2-density result for event structures}
\label{subsec:2density}

\newcommand{\bE}{\mathbf{E}}
\newcommand{\bC}{\mathbf{C}}
\newcommand{\bD}{\mathbf{D}}

We prove a density result for the embedding of event structures into event structures with symmetry. This is a 2-categorical (i.e. $\Cat$-enriched) form of density: each event structure with symmetry is a canonical 2-colimit of  ordinary event structures. Another way of stating this is as follows:

\newcommand{\op}{\mathrm{op}}

\begin{definition}[e.g. [17]]
For locally small 2-categories $\bD$ and $\bE$, a $2$-functor $J : \bD \rightarrow \bE$ is \emph{2-dense} if the 2-functor $\tilde{J} : \bE \to [\bD^\op, \Cat] : e \mapsto \bE(J-, e)$ is locally an isomorphism of categories, where $[\bD^\op, \Cat]$ is the 2-category of 2-functors, strict natural transformations, and modifications. 
\end{definition}

\begin{theorem}
The embedding 2-functor $J : \ES \to \SES$ is 2-dense.
\end{theorem}

(We emphasize that 2-density is strictly weaker than density in the 1-categorical sense. In particular the underlying 1-functor $J : \ES \to \SES$ is not dense.)
\begin{proof}
For event structures with symmetry $A$ and $B$, we show that the functor $\tilde{J}_{A ,B} : \SES(A, B) \to \Cat^{\ES^\op}(\SES(J-, A), \SES(J-, B))$ is an isomorphism of categories by constructing an inverse. Write $A = (|A|, \mathbb{S}_A)$ and $B = (|B|, \mathbb{S}_B)$. 

Observe that, for a natural transformation $\alpha : \SES(J-, A) \to \SES(J-, B)$, applying component $\alpha_{|A|}$ to the `identity' function $\id_{|A|} : J|A| \to A$ gives a map of event structures with symmetry $\alpha_{|A|}(\id_{|A|}) : J|A| \to B$. We claim that $\alpha_{|A|}(\id_{|A|})$ preserves the symmetry in $A$. For any bijection $\theta: x \cong y$ in $\mathbb{S}_A$, we have $(\id_{|A|} \circ J\iota_y \circ J\theta) \sim (\id_{|A|} \circ J\iota_x)$, writing $\iota_x : x \to |A|$ and $\iota_y : y \to |A|$ for the inclusion maps (and recalling that $\theta \in \ES(x, y)$). Since $\alpha$ is a natural transformation and each component preserves the symmetry relation on maps,  $\alpha_{|A|}(\id_{|A|}) \circ J{\iota_x} = \alpha_x (\id_{|A|} \circ J{\iota_x}) \sim \alpha_x(\id_{|A|} \circ J{\iota_y} \circ J\theta) = \alpha_{|A|}(\id_{|A|}) \circ J\iota_y \circ J\theta.$ In other words, $\alpha_{|A|}(\id_{|A|}) \theta \in \mathbb{S}_B$, and we have proven the claim. Therefore, we take $\tilde J_{A, B}^{-1}(\alpha) = \alpha_{|A|}(\id_{|A|})$. For functoriality of $J_{A, B}^{-1}$, note that any modification $\alpha \to \beta$ must be unique and gives $\alpha_{|A|}(\id_{|A|})\sim \beta_{|A|}(\id_{|A|})$. We omit the verification that $\tilde J_{A, B}^{-1}$ and $\tilde J_{A, B}$ are inverses. 
\end{proof}

\begin{thebibliography}{99}
\bibitem[1]{source1} Nathanael Arkor and Dylan McDermott. The formal theory of relative monads. Journal of Pure and Applied Algebra, 228(9):107676, September 2024. URL: https://doi.org/10.1016/j.jpaa.2024.107676 .
\bibitem[2]{source2} John C Baez, Fabrizio Genovese, Jade Master, and Michael Shulman. Categories of nets. In 2021 36th Annual ACM/IEEE Symposium on Logic in Computer Science (LICS), pages 1-13. IEEE, 2021. URL: https://doi.org/10.1109/LICS52264.2021.9470566 .
\bibitem[3]{source3} John C Baez and Jade Master. Open petri nets. Mathematical Structures in Computer Science, 30(3):314-341, 2020. URL: https://doi.org/10.1017/S0960129520000043 .
\bibitem[4]{source4} Kris Brown, Kevin Carlson, Owen Lynch, and Evan Patterson. Catcolab software. URL: https://catcolab.org/help/ .
\bibitem[5]{source5} Roberto Bruni, José Meseguer, Ugo Giovanni Erasmo Montanari, and Vladimiro Sassone. Functorial models for petri nets. Information and Computation, 170(2):207-236, 2001. URL: https://doi.org/10.1006/inco.2001.3050 .
\bibitem[6]{source6} Benjamin Merlin Bumpus, Sophie Libkind, Jordy Lopez Garcia, Layla Sorkatti, and Samuel Tenka. Additive invariants of open petri nets. Compositionality, 7, 2025.
\bibitem[7]{source7} Aurelio Carboni, Stefano Kasangian, and Ross Street. Bicategories of spans and relations. Journal of pure and applied algebra, 33(3):259-267, 1984.
\bibitem[8]{source8} Simon Castellan and Pierre Clairambault. The geometry of causality: Multi-token geometry of interaction and its causal unfolding. Proc. ACM Program. Lang., 7(POPL):689-717, 2023. URL: https://doi.org/10.1145/3571217 .
\bibitem[9]{source9} Simon Castellan, Pierre Clairambault, Silvain Rideau, and Glynn Winskel. Games and strategies as event structures. Logical Methods in Computer Science, Volume 13, Issue 3, September 2017. URL: https://doi.org/10.23638/lmcs-13(3:35)2017 .
\bibitem[10]{source10} Simon Castellan, Pierre Clairambault, and Glynn Winskel. The parallel intensionally fully abstract games model of pcf. In 2015 30th Annual ACM/IEEE Symposium on Logic in Computer Science, pages 232-243. IEEE, 2015. URL: https://doi.org/10.1109/LICS.2015.31 .
\bibitem[11]{source11} David Dalrymple. Safeguarded ai: Constructing safety by design. Advanced Research and Invention Agency, 2024. URL: https://aria.org.uk/media/3nhijno4/aria-safeguarded-ai-programme-thesis-v1.pdf .
\bibitem[12]{source12} Elena Di Lavore, Wilmer Leal, and Valeria de Paiva. Dialectica petri nets. Fundamenta Informaticae, 194, 2025. URL: https://doi.org/10.46298/fi.13125 .
\bibitem[13]{source13} Uli Fahrenberg, Christian Johansen, Christopher Trotter, and Krzysztof Ziemianski. Sculptures in concurrency. Log. Methods Comput. Sci., 17(2), April 2021. URL: https://doi.org/10.23638/lmcs-17(2:5)2021 .
\bibitem[14]{source14} Eric Goubault and Samuel Mimram. Formal relationships between geometrical and classical models for concurrency. In Lisbeth Fajstrup, Eric Goubault, and Martin Raussen, editors, Proceedings of the workshop on Geometric and Topological Methods in Computer Science, GETCO 2010, Aalborg, Denmark, January 11-15, 2010, volume 283 of Electronic Notes in Theoretical Computer Science, pages 77-109. Elsevier, June 2010. URL: https://doi.org/10.1016/j.entcs.2012.05.007 .
\bibitem[15]{source15} Jonathan Hayman and Glynn Winskel. Symmetry in petri nets. Unpublished manuscript, January 2008.
\bibitem[16]{source16} Jonathan Hayman and Glynn Winskel. The unfolding of general Petri nets. In Ramesh Hariharan, Madhavan Mukund, and V Vinay, editors, IARCS Annual Conference on Foundations of Software Technology and Theoretical Computer Science, volume 2 of Leibniz International Proceedings in Informatics (LIPIcs), pages 223-234, Dagstuhl, Germany, 2008. Schloss Dagstuhl - Leibniz-Zentrum für Informatik. URL: https://doi.org/10.4230/LIPIcs.FSTTCS.2008.1755 .
\bibitem[17]{source17} G. M. Kelly. Basic Concepts of Enriched Category Theory, volume 64 of London Mathematical Society Lecture Note Series. Cambridge University Press, 1982. Reprinted in: Reprints in Theory and Applications of Categories, No. 10 (2005), pp. 1-136.
\bibitem[18]{source18} Joachim Kock. Whole-grain petri nets and processes. J. ACM, 70(1), December 2022. URL: https://doi.org/10.1145/3559103 .
\bibitem[19]{source19} Serge Lechenne, Clovis Eberhart, and Ichiro Hasuo. A compositional framework for petri nets. In Barbara König and Henning Urbat, editors, Coalgebraic Methods in Computer Science - 17th IFIP WG 1.3 International Workshop, CMCS 2024, Colocated with ETAPS 2024, Luxembourg City, Luxembourg, April 6-7, 2024, Proceedings, volume 14617 of Lecture Notes in Computer Science, pages 174-193, Berlin, Heidelberg, 2024. Springer. URL: https://doi.org/10.1007/978-3-031-66438-0\_9 .
\bibitem[20]{source20} Jade Master and Joe Moeller. Colored petri nets are monoidal double functors. arXiv preprint, 2025. URL: https://arxiv.org/abs/2510.01946 .
\bibitem[21]{source21} Kenneth L. McMillan. Using unfoldings to avoid the state explosion problem in the verification of asynchronous circuits. In Gregor von Bochmann and David K. Probst, editors, Computer Aided Verification, Fourth International Workshop, CAV '92, Montreal, Canada, June 29 - July 1, 1992, Proceedings, volume 663 of Lecture Notes in Computer Science, pages 164-177, Berlin, Heidelberg, 1992. Springer. URL: https://doi.org/10.1007/3-540-56496-9\_14 .
\bibitem[22]{source22} José Meseguer and Ugo Montanari. Petri nets are monoids. Information and computation, 88(2):105-155, 1990. URL: https://doi.org/10.1016/0890-5401(90)90013-8 .
\bibitem[23]{source23} José Meseguer, Ugo Montanari, and Vladimiro Sassone. On the semantics of place/transition petri nets. Mathematical Structures in Computer Science, 7(4):359-397, 1997. URL: https://doi.org/10.1017/S0960129597002314 .
\bibitem[24]{source24} T. Murata. Petri nets: Properties, analysis and applications. Proceedings of the IEEE, 77(4):541-580, 1989. URL: https://doi.org/10.1109/5.24143 .
\bibitem[25]{source25} Mogens Nielsen, Gordon D. Plotkin, and Glynn Winskel. Petri nets, event structures and domains, part I. Theor. Comput. Sci., 13:85-108, 1981. URL: https://doi.org/10.1016/0304-3975(81)90112-2 .
\bibitem[26]{source26} Wolfgang Reisig. Petri Nets: An Introduction, volume 4 of EATCS Monographs on Theoretical Computer Science. Springer, 1985. URL: https://doi.org/10.1007/978-3-642-69968-9 .
\bibitem[27]{source27} Sam Staton and Glynn Winskel. On the expressivity of symmetry in event structures. In 2010 25th Annual IEEE Symposium on Logic in Computer Science, pages 392-401, 2010. URL: https://doi.org/10.1109/LICS.2010.37 .
\bibitem[28]{source28} Rob J van Glabbeek and Gordon D Plotkin. Configuration structures, event structures and petri nets. Theoretical Computer Science, 410(41):4111-4159, 2009. URL: https://doi.org/10.1016/j.tcs.2009.06.014 .
\bibitem[29]{source29} Glynn Winskel. A new definition of morphism on petri nets. In Annual Symposium on Theoretical Aspects of Computer Science, pages 140-150. Springer, 1984. URL: https://doi.org/10.1007/3-540-12920-0\_13 .
\bibitem[30]{source30} Glynn Winskel. Event structures. In Wilfried Brauer, Wolfgang Reisig, and Grzegorz Rozenberg, editors, Petri Nets: Central Models and Their Properties, Advances in Petri Nets 1986, Part II, Proceedings of an Advanced Course, Bad Honnef, Germany, 8-19 September 1986, volume 255 of Lecture Notes in Computer Science, pages 325-392, Berlin, Heidelberg, 1986. Springer. URL: https://doi.org/10.1007/3-540-17906-2\_31 .
\bibitem[31]{source31} Glynn Winskel. Petri nets, algebras, morphisms, and compositionality. Information and Computation, 72(3):197-238, 1987. URL: https://doi.org/10.1016/0890-5401(87)90032-0 .
\bibitem[32]{source32} Glynn Winskel. Event structures with symmetry. In Luca Cardelli, Marcelo Fiore, and Glynn Winskel, editors, Computation, Meaning, and Logic: Articles dedicated to Gordon Plotkin, volume 172 of Electronic Notes in Theoretical Computer Science, pages 611-652. Elsevier, 2007. Computation, Meaning, and Logic: Articles dedicated to Gordon Plotkin. URL: https://doi.org/10.1016/j.entcs.2007.02.022 .
\end{thebibliography}
\end{document}
